\documentclass[12pt,a4paper]{article}
\usepackage[utf8]{inputenc}
\usepackage[vietnamese]{babel}
\usepackage[margin=2cm]{geometry}
\usepackage{tcolorbox}
\usepackage{fancyhdr}

\pagestyle{fancy}
\fancyhf{}
\lhead{\textbf{VIỆN KHOA HỌC HÌNH SỰ}}
\rhead{\textbf{KẾT LUẬN GIÁM ĐỊNH}}
\cfoot{Trang \thepage}

\begin{document}

\begin{center}
    \textbf{\Large CỘNG HÒA XÃ HỘI CHỦ NGHĨA VIỆT NAM}\\
    \textbf{\large Độc lập - Tự do - Hạnh phúc}\\
    \rule{8cm}{0.5pt}\\[0.5cm]
    
    \textbf{\LARGE KẾT QUẢ GIÁM ĐỊNH PHÁP Y CHỮ KÝ \& MỰC}\\
    \textbf{Số hiệu:} \texttt{KLGD-2026/88}
\end{center}

\vspace{0.3cm}

\begin{tcolorbox}[colback=gray!10!white,colframe=black!75!black,title=\textbf{THÔNG TIN GIÁM ĐỊNH}]
\begin{itemize}[leftmargin=*]
    \item \textbf{Cơ quan giám định:} Viện Khoa học Hình sự
    \item \textbf{Yêu cầu từ:} Luật sư Nguyễn Văn Minh (Đại diện cho Trần Ngọc Mai)
    \item \textbf{Ngày ra kết luận:} 25/07/2026
\end{itemize}
\end{tcolorbox}

\section*{I. NỘI DUNG GIÁM ĐỊNH}
\begin{enumerate}
    \item \textbf{Mẫu giám định:} Bản gốc di chúc viết tay đứng tên Nguyễn Văn Thành do Trần Ngọc Mai nộp.
    \item \textbf{Phương pháp:} Giám định áp quang quang phổ \& sắc ký mực.
\end{enumerate}

\section*{II. KẾT LUẬN GIÁM ĐỊNH}
\begin{itemize}
    \item Phần chữ ký \textit{"Nguyễn Văn Thành"} là chữ ký thật của cụ Thành.
    \item Tuy nhiên, đoạn văn bản mô tả \textit{"thuộc Mã thửa địa chính khóa 2021-BS14 cho cháu nội Nguyễn Văn Khang..."} được viết bằng \textbf{chất liệu mực bi hóa dầu sản xuất năm 2024}, viết chèn đè lên khoảng trống tờ giấy sau khi cụ Thành đã ký tên.
    \item \textbf{Kết luận chính thức:} Tờ di chúc bị chỉnh sửa và làm giả nội dung nhằm thâu tóm tài sản.
\end{itemize}

\vspace{1cm}

\begin{tcolorbox}[colback=blue!5!white,colframe=blue!75!black,title=\textbf{Ý NGHĨA PHÁP LÝ}]
Kết luận này chứng minh Trần Ngọc Mai đã chủ động mang di chúc đi giám định hợp pháp và không có động cơ giết người tại hiện trường.
\end{tcolorbox}

\end{document}
